\chapter{Analysis \& Discussion} \label{ch:analysis_and_discussion}
%Should include a reiteration of the experiments, and their outcome.  Together with a description (discussion).  Preamble should include a reminder of the aims and objectives together with a list of experiments to achieve these.  Should include many charts and other visualization with appropriate descriptions.\footnote{Another footnote example.}  

%TODO Given that both games in this section feature eponymous protagonists, I will be distinguishing between the two by writing the names of the games in italics and those of the characters in a regular font.

\section{Alum}

\subsection{Game Overview}

\textit{Alum} \parencite{orsie_alum_2006} is a point-and-click adventure game heavily inspired by older titles developed by LucasArts and Sierra Entertainment. Like those games, \textit{Alum} features an overarching scripted narrative which players must progress by interacting with the environment. They can do so through an avatar capable of performing different actions (i.e., ``verbs'') and an inventory, which they can use to store and combine items. Occasionally, \textit{Alum} adds some variety by requiring the player to complete action sequences, such as a sniper mini-game. Furthermore, the also game allows the player to diverge from the main sequence of scripted events in three occasions. However, in each case, doing so plays a brief narrative sequence, which leads directly to a game over screen.

Starting a new game introduces us to Alum Descry, a resident of Kozmos city, located within the land of Tide. Just like many other inhabitants, Alum's wife, Esther, is afflicted by a disease known as ``The Vague'', which causes her to perceive life as meaningless and withdraw from others, mostly remaining silent. 

%Determined to find a cure, Alum tries to arrange a meeting with members of a resistance which claim to have a cure. Once they do meet however, things go awry and Alum ends up having to escape the city.


 
\subsection{Key Observations} \label{sec:alum_key_observations}

\section{Captain Bible in Dome of Darkness}

\subsection{Game Overview}

\subsection{Key Observations} \label{sec:captain_key_observations}

\section{Summary}
%\blindtext\enlargethispage{\baselineskip} % so you do not get a single line in another page